\section{Transitive permutation quotients and imprimitive Heisenberg carriers}
\label{sec:heisenberg-permutation-carrier}

\paragraph{Author freeze and provenance.}
D-TD40, reserved question Q2ev, is a root derivation.
The initial coarsest-block proposal is strengthened here to every
transitive permutation image by a finite induction on block degree.
The actual coordinate-kernel mechanism is the accepted D-TD17
argument; the new induction retains the quotient through every
smaller block action.  The published minimum-degree input is stated
with its exact locator below.  The complete author written-proof
and source-interface readings pass; independent whole-proof review
of this new section also passes.

\subsection*{D-TD40(a). A support alternative for an actual transitive quotient}

Write $\operatorname{mov}_X(g)$ for the number of elements of $X$
moved by a permutation $g$.  Let $n\ge10^6$ and let
\[
 G\le\operatorname{Sym}(\Omega)\text{ be transitive},\qquad
 t=|\Omega|,\qquad
 \varphi:G\twoheadrightarrow T\in\{A_n,S_n\}.
 \tag{TPC.1}
\]
The map is an actual surjection, with its specified natural action
on $n$ letters.  Then $t\ge n$.  Moreover at least one of the
following two alternatives holds for the whole pair $(G,\varphi)$:
\[
 \begin{split}
 &\text{there is }u\mid t,\quad n\le u\le t,\quad
 \operatorname{mov}_{\Omega}(g)
       \ge\frac{t}{u}\frac{\sqrt u-1}{2}
       \quad\text{whenever }\varphi(g)\ne1; \qquad\text{or}\\
 &n\mid t,\qquad
 \operatorname{mov}_{\Omega}(g)
       \ge\frac{t}{n}\operatorname{mov}_{[n]}(\varphi(g))
       \quad\text{for every }g\in G.
 \end{split}
 \tag{TPC.2}
\]
The integer $u$ and the choice of alternative do not depend on $g$.

\begin{proof}
Since $|T|\ge n!/2$ and $|G|\le t!$, the inequality $t<n$
would imply $|G|\le(n-1)!<n!/2$, a contradiction.

We use induction on $t$.  Choose a maximal proper block system
$\mathcal B$ for the transitive action of $G$ on $\Omega$.
Its blocks have a common size $b$, their number is $v=t/b>1$,
and the induced permutation image $\Pi\le S_v$ is primitive.
For a primitive $G$, this means the singleton block system.
Indeed a proper nontrivial block system for $\Pi$ would merge
the original blocks to a strictly coarser proper block system.
Let $K$ be the actual kernel of $G\to\Pi$.

Suppose first that $\varphi(K)=1$.  Then there is an actual
surjection $\bar\varphi:\Pi\twoheadrightarrow T$, and the same
order comparison gives $v\ge n$.  If $\Pi$ does not contain
$A_v$, the published theorem of Praeger--Shalev gives
\[
 \min_{1\ne\pi\in\Pi}\operatorname{mov}_{\mathcal B}(\pi)
                       \ge\frac{\sqrt v-1}{2}.
 \tag{TPC.3}
\]
Here the precise source is \emph{Bounds on finite quasiprimitive
permutation groups}, J. Austral. Math. Soc. \textbf{71} (2001),
243--258, Theorem~7.2, pp.~255--256.\cite{TD:PraegerShalev2001}
A primitive group is quasiprimitive: the orbits of a normal subgroup
form a block system, and singleton orbits would make that subgroup
trivial in the faithful permutation action.  Thus the source applies
to this actual primitive image.  Its supplied proof and surrounding
discussion cite the earlier Babai and Manning primitive bounds;
no new re-proof of those upstream results is asserted here.
Whenever $\varphi(g)\ne1$, its image in $\Pi$ is nonidentity.
Every moved block contributes all $b$ of its points to the support
of $g$.  This proves the first alternative of (TPC.2) with $u=v$.

If $\Pi$ contains $A_v$, it is $A_v$ or $S_v$ in its natural
action.  Since $v\ge n\ge10^6$, the elementary normal-subgroup
lists show that a nonabelian quotient of $\Pi$ is $\Pi$ itself:
the proper nontrivial normal subgroup of $S_v$ is $A_v$, whose
quotient has order two, and $A_v$ is simple.
Thus $\bar\varphi$ is an isomorphism.  Its derived subgroup
identifies $A_v$ with $A_n$, so the strictly increasing factorial
orders give $v=n$.  The actual natural quotient is identified
up to a relabeling of the $n$ letters.  For $A_n$, this is the
accepted $\operatorname{Aut}(A_n)=S_n$ input used in D-TD35.
For $S_n$, restrict the automorphism to $A_n$ and remove that
inner conjugation; the resulting automorphism fixes $A_n$
pointwise, and for every $x\in S_n$ its value differs from $x$
by an element of $C_{S_n}(A_n)=1$.  Hence it too is inner.
Relabeling preserves support.  The moved blocks of $g$ therefore
number $\operatorname{mov}_{[n]}(\varphi(g))$, proving the
second alternative of (TPC.2).

It remains to consider $\varphi(K)\ne1$.  Put $T_0=\varphi(K)$.
Normality in $T$ gives
\[
                         A_n\le T_0\le T.
 \tag{TPC.4}
\]
In particular $T_0$ is $A_n$ or $S_n$.  For every block
$B\in\mathcal B$, let $L_B=\ker(K\to\operatorname{Sym}(B))$.
These subgroups are normal in $K$ and have trivial common
intersection, because the action of $G$ on $\Omega$ is faithful.
The actual homomorphism $K\to T_0$ kills at least one $L_B$.
To prove this explicitly, if every $\varphi(L_B)$ were nontrivial,
each would contain $A_n$.  For successive intersections $I_j$,
\[
 [I_{j-1},L_{B_j}]\le I_j,
 \qquad
 \varphi([I_{j-1},L_{B_j}])
   =[\varphi(I_{j-1}),\varphi(L_{B_j})]\supseteq A_n.
 \tag{TPC.5}
\]
Induction would make the trivial final intersection have image
containing $A_n$, a contradiction.  Conjugation by $G$ permutes
the $L_B$ transitively and preserves $\ker\varphi$.  Consequently
\[
                   L_B\le\ker\varphi
                         \quad\text{for every }B\in\mathcal B.
 \tag{TPC.6}
\]
No transitivity of $\ker\varphi$ is used.

Fix one block $B$ and its setwise stabilizer $G_B$.
The induced image $\Gamma_B\le\operatorname{Sym}(B)$ is
transitive: if $x,y\in B$, any $g\in G$ with $gx=y$ has
$gB\cap B\ne\varnothing$ and therefore $gB=B$.
If $h\in G_B$ acts identically on $B$, then for every $k\in K$
the commutator $[h,k]$ lies in $L_B$.  By (TPC.6),
$\varphi(h)$ centralizes $\varphi(K)\supseteq A_n$.
The centralizer of $A_n$ in $S_n$ is trivial.  One elementary
verification is that a commuting permutation sends a letter to
a letter fixed by its natural $A_{n-1}$ stabilizer; for $n\ge5$
that stabilizer fixes only the original letter.  Thus
$\varphi(h)=1$.  The quotient therefore descends to an actual
surjection
\[
 \varphi_B:\Gamma_B\twoheadrightarrow T_B:=\varphi(G_B),
                  \qquad A_n\le T_B\le T.
 \tag{TPC.7}
\]
This conclusion is proved for the actual restriction image;
it is not inferred from a composition factor.

The degree $b$ is strictly smaller than $t$, so induction applies
to $(\Gamma_B,\varphi_B)$, with the same natural $n$.
In particular $b\ge n$.  Write its conclusion in the form
\[
 \operatorname{mov}_{B}(h|_B)\ge b\,\eta(\varphi(h)),
 \qquad 0\le\eta(y)\le1,
 \tag{TPC.8}
\]
where either $u\mid b$, $u\ge n$ and
$\eta(y)=(\sqrt u-1)/(2u)$ for $y\ne1$, with $\eta(1)=0$,
or $u=n\mid b$ and $\eta(y)=\operatorname{mov}_{[n]}(y)/n$.
The function is invariant under conjugation in $S_n$.
Transporting $B$ by elements of $G$ gives the same inequality,
with the same $u$ and same function $\eta$, on every block
fixed setwise by any given $g\in G$.  More precisely, conjugating
such a restriction back to $B$ conjugates its natural image by
the corresponding element of $T$, and does not change $\eta$.
Every block moved by $g$ contributes all $b$ points, which is
at least $b\eta(\varphi(g))$.  Summing over the $t/b$ blocks gives
\[
             \operatorname{mov}_{\Omega}(g)
                            \ge t\,\eta(\varphi(g)).
 \tag{TPC.9}
\]
Since $u\mid b\mid t$, this is exactly (TPC.2).  The induction
terminates by strict degree decrease.  It supplies one alternative
valid for all elements, rather than choosing a branch for each $g$.
\end{proof}

\subsection*{D-TD40(b). Exclusion of an entire block-permutation carrier}

Use the actual stable whole Heisenberg setup of D-TD34, with
$Q\simeq\mathbf F_p^{2m}$, actual irreducible faithful defining
image $H$, actual surjection $\psi:H\twoheadrightarrow T$
for $T=A_n$ or $S_n$, and $n\ge10^6$.
Let $F=\operatorname{End}_{\mathbf F_pH}(Q)=\mathbf F_{p^f}$,
$a=\dim_FQ$ and $M=Q\otimes_F\overline F$.
This absolute model is irreducible by the accepted finite-algebra
argument; suppose it admits an actual $H$-stable imprimitivity system
\[
 M=\bigoplus_{\omega\in\Omega}M_\omega,\qquad
 \dim_{\overline F}M_\omega=b,\qquad t=|\Omega|>1,
 \qquad a=bt.
 \tag{TPC.10}
\]
Irreducibility makes the block action transitive.  Suppose the
actual $\psi$ kills the kernel of that permutation action.
Then no actual sphere source with this entire Heisenberg factor,
arbitrary entire external complex partner $R$, and moving-rank
sum $2e+6$ can exist.

\begin{proof}
Let $\Pi\le S_t$ be the actual transitive block image.
The hypothesis supplies a genuine surjection $\Pi\to T$.
Apply (TPC.2) to this pair, retaining its single choice of alternative.

For any $h\in H$, let $\pi$ be its permutation of the summands.
On an $\ell$-cycle of summands, an eigenvector at any nonzero
target is determined by its component in the first summand.
Thus that cycle has kernel dimension at most $b$ and rank at
least $(\ell-1)b$.  This follows from the successive equations
around the cycle, in every characteristic; diagonalizability is
not used.  The target zero gives full rank because $h$ is invertible.
Writing $\operatorname{ind}_\Omega(\pi)=t-\#\{\pi\text{-cycles}\}$,
we obtain
\[
 \operatorname{rank}_{\mathbf F_p}(1-h)
  =f\operatorname{rank}_{\overline F}(1-h|_M)
  \ge fb\operatorname{ind}_{\Omega}(\pi)
  \ge\frac{fb}{2}\operatorname{mov}_{\Omega}(\pi).
 \tag{TPC.11}
\]
At every scalar target over $\overline F$, the unscaled rank is
at least $b\operatorname{ind}_{\Omega}(\pi)$.  The factor $f$
in (TPC.11) concerns the original target-one prime-field rank.
Individual blocks need not descend to $F$;
rank extension and restriction are applied only to the original
$F$-matrix of $h$.

In the first alternative, (TPC.11) yields for every $\psi(h)\ne1$
\[
 \operatorname{rank}_{\mathbf F_p}(1-h)
       \ge\frac{fa}{u}\frac{\sqrt u-1}{4}
       \ge\frac{999}{4}>13,
 \qquad u\ge n\ge10^6,\qquad u\mid t\le a.
 \tag{TPC.12}
\]
The exact D-TD05 Heisenberg trace formula therefore gives normalized
trace at most $c=2^{-13/2}<1/90$ for every original coefficient
element with nonidentity natural image.  The entire partner is
retained.  For an original element of order $L$ whose natural
image has order $o$, exactly $L/o$ of its powers have trivial
natural image; their traces have normalized absolute value at
most one.  Fourier inversion over all $L$ powers gives moving
rank at least $e(1-c)(1-1/o)$.
The actual connected $T$ quotient of the original sphere cover
has genus at least two and orbifold sum at least $85/42$, by the
actual-cover argument of D-TD02.  The original budget consequently
gives $e(1-85c)\le252$, and hence $e<4536$.
But $a\ge t\ge n$ and $\dim W=p^{fa/2}$ give
$e\ge2^{n/2}>4536$, a contradiction.

In the second alternative put $k=\operatorname{mov}_{[n]}(\psi(h))$.
The support and rank bounds give
\[
 \operatorname{rank}_{\mathbf F_p}(1-h)
                         \ge\frac{fa}{2n}k\ge\frac{k}{2}.
 \tag{TPC.13}
\]
For a nonidentity natural image this is at least one.  The normalized
Heisenberg trace is therefore at most $1/\sqrt2<7/8$.
If $k\le n/16$, the exact inequality
$2\cdot15^6>16^6$ gives the stronger upper bound
$1/\sqrt2<(15/16)^3\le((n-k)/n)^3$.
If $k>n/16$, (TPC.13) gives trace less than
$2^{-n/64}<n^{-3}$; the accepted elementary dyadic bound in
D-TD22 is already stronger, namely $2^{1-n/64}<n^{-3}$ for
$n\ge4096$.  Also $e\ge2^{a/2}\ge2^{n/2}>n^3$, by that same
proved degree estimate.  Thus every element of the actual whole
coefficient representation satisfies both cubic and uniform
$7/8$ trace hypotheses of HD-FT27.

For clarity, the whole-partner comparison uses finite projective
factor images and their finite special-linear lifts, followed by
averaged positive Hermitian forms.  Every original tensor operator
is a scalar times two unitary factor representatives; finite order
makes that scalar have absolute value one.  The whole normalized
trace is bounded by the $W$ trace for every element and every full
original power.  Natural-kernel powers retain the bound one.
No quotient order replaces $L$, and no cover of a block stabilizer
or intermediate permutation constituent is assumed.

If the coefficient representation is not faithful, pass to the
actual connected cover $C/\ker V$.  A tensor operator equal to
the identity has scalar factors, so $\ker V$ kills the projective
$W$ image and hence $H$ and $T$.  The quotient retains the actual
$T$ map, all character bounds and the same moving-rank sum;
identity coefficient inertia images contribute zero.
The faithful-source formulation of HD-FT27 now excludes the budget.
\end{proof}

\paragraph{Exact remaining boundary.}
No primitivity of the block permutation image is assumed in part(b).
The initial block system can be finest, coarsest, or any intermediate
system.  The theorem requires that its actual block-permutation kernel
be killed by $\psi$.  If it has nontrivial image, normality forces
that image to contain $A_n$, and the separate block-diagonal
argument is required.  The result supplies no target polarization,
new detecting cycle, or bound for groups without this actual quotient.
